\chapter{Algebraic foundations and bar constructions}\label{ch:algebraic-foundations}

\section{Classical Koszul duality}
\label{sec:classical-koszul-foundation}

Every duality in this monograph---Verdier on configuration spaces,
complementarity of quantum corrections, the Koszul involution
$k \leftrightarrow -k-2h^\vee$ on affine levels---is a shadow of a
single algebraic mechanism: the bar-cobar adjunction for quadratic
algebras.  This chapter develops that mechanism in its purest form,
before geometry and analysis enter.

The reader should regard the constructions here as the genus-$0$ seed
of the full modular theory.  Bar and cobar are adjoint functors on
dg-categories; the Koszul property---diagonal Ext vanishing---controls
when the adjunction is an equivalence; and the combinatorics of the
tensor coalgebra $T^c(s^{-1}\bar{A})$ already encodes, in skeleton,
the boundary stratification of the Fulton--MacPherson
compactification.  When configuration spaces on curves replace tensor
products in Chapter~\ref{chap:config-spaces}, and residues on
logarithmic forms replace the alternating-sum differential in
Chapter~\ref{chap:bar-cobar}, the formal properties established here
--- $d^2 = 0$ from quadraticity, the Koszul criterion, the adjunction
--- will persist as geometric theorems rather than algebraic identities.
Without this template, the geometric construction would have no
structural anchor.

\subsection{Quadratic algebras and Koszul duality}

\begin{definition}[Quadratic algebra]
\label{def:quadratic-algebra}
\index{quadratic algebra|textbf}
A graded algebra $A = T(V)/I$ is \emph{quadratic} if:
\begin{enumerate}
\item $V$ is a graded vector space (generators)
\item $I \subset V \otimes V$ is a subspace of relations in degree 2
\item The defining ideal is $(I)$ generated by $I$
\end{enumerate}
We write $A = A(V, R)$ where $R \subset V \otimes V$ are the relations.
\end{definition}

\begin{example}[Prototypical examples]
\label{ex:classical-quadratic-algebras}
\index{symmetric algebra}
\index{exterior algebra}
\index{universal enveloping algebra}
\index{tensor algebra}
\begin{enumerate}
\item \emph{Commutative algebra} $\text{Sym}(V)$:
\begin{align}
\text{Generators:} &\quad V \\
\text{Relations:} &\quad R_{\text{Com}} = \{v \otimes w - w \otimes v : v,w \in V\} \subset V \otimes V
\end{align}

\item \emph{Exterior algebra} $\Lambda(V)$:
\begin{align}
\text{Generators:} &\quad V \\
\text{Relations:} &\quad R_{\Lambda} = \{v \otimes w + w \otimes v : v,w \in V\} \subset V \otimes V
\end{align}

\item \emph{Universal enveloping} $U(\mathfrak{g})$ for Lie algebra $\mathfrak{g}$:
\begin{align}
\text{Generators:} &\quad \mathfrak{g} \\
\text{Relations:} &\quad R_{\mathfrak{g}} = \{x \otimes y - y \otimes x - [x,y] : x,y \in \mathfrak{g}\}
\end{align}
\end{enumerate}
\end{example}

\subsection{The Koszul dual coalgebra}

\begin{construction}[Quadratic dual]
\label{const:quadratic-dual}
\index{Koszul dual!quadratic|textbf}
Given a quadratic algebra $A = A(V,R)$, define its \emph{quadratic dual} $A^! = A(V^*, R^\perp)$ by:
\begin{align}
\text{Generators:} &\quad V^* \quad \text{(dual space)} \\
\text{Relations:} &\quad R^\perp = \{r \in V^* \otimes V^* : \langle r, s \rangle = 0 \text{ for all } s \in R\}
\end{align}
where the pairing is:
\begin{equation}
\langle \alpha \otimes \beta, v \otimes w \rangle = \langle \alpha, v \rangle \langle \beta, w \rangle
\end{equation}
\end{construction}

\begin{remark}[Orthogonality principle]
The spaces $R$ and $R^\perp$ are orthogonal complements in $V \otimes V$ and $V^* \otimes V^*$ respectively; this orthogonality is the source of duality.
\end{remark}

\subsection{Koszul pairs: precise definition}

\begin{definition}[Koszul pair]
\label{def:koszul-pair-classical}
\index{Koszul pair!classical|textbf}
A pair of quadratic algebras $(A_1, A_2)$ is a \emph{Koszul pair} if:
\begin{enumerate}
\item $\bar{B}(A_1) \simeq A_2^!$ (as coalgebras)
\item $\bar{B}(A_2) \simeq A_1^!$ (as coalgebras)
\item $\Omega(\bar{B}(A_1)) \simeq A_1$ (cobar inverts bar)
\item $\Omega(\bar{B}(A_2)) \simeq A_2$ (cobar inverts bar)
\end{enumerate}
\end{definition}

\begin{remark}[Two phenomena]
\label{rem:two-phenomena}
Conditions (1-2) establish \emph{Koszul duality}: $A_1$ and $A_2$ encode dual coalgebraic information.

Conditions (3-4) establish \emph{bar-cobar inversion}: the composite $\Omega \circ \bar{B}$ is homotopy equivalent to the identity.

These are distinct mathematical phenomena:
\begin{itemize}
\item $\bar{B}(A_1) \simeq A_2^!$ means that the bar of $A_1$ produces the dual coalgebra to $A_2$;
\item $\Omega(\bar{B}(A_1)) \simeq A_1$ means that cobar reconstructs $A_1$ from its bar coalgebra;
\item together, $A_1$ and $A_2$ are Koszul dual, with bar-cobar mediating the duality.
\end{itemize}
\end{remark}

\subsection{Classical examples revisited}

\begin{theorem}[Classical Koszul pairs; \ClaimStatusProvedElsewhere]
\label{thm:quadratic-koszul}
\index{Koszul duality!classical}
\index{Chevalley--Eilenberg complex}
The following are Koszul pairs in the sense of Definition \ref{def:koszul-pair-classical}:
\begin{enumerate}
\item $(\text{Sym}(V), \Lambda(V^*))$ — commutative and exterior algebras
\item $(U(\mathfrak{g}), C^*_{\text{CE}}(\mathfrak{g}))$ — universal enveloping and Chevalley--Eilenberg cochains
\item $(T(V), T^c(V^*))$ — tensor algebra and tensor coalgebra
\end{enumerate}
Each pair satisfies all four conditions of Definition \ref{def:koszul-pair-classical}.
\end{theorem}

\begin{remark}[Provenance and citation]
This theorem is imported and treated as \ClaimStatusProvedElsewhere. Classical
Koszul duality for these standard quadratic examples is foundational material in
\cite{GK94,LV}.
\index{Loday--Vallette}
\end{remark}

We now view early examples of the chiral enhancement of this classical structure.


\emph{Free fermions.}

Let $\mathcal{F}$ be the free fermion chiral algebra with generator $\psi(z)$ and OPE:
$$\psi(z)\psi(w) \sim \frac{1}{z-w}$$

Computation shows:
$$\Omega(\bar{B}(\mathcal{F})) \simeq \mathcal{F}$$


\emph{Heisenberg algebra.}

Let $\mathcal{H}_k$ be the Heisenberg chiral algebra with generator $\alpha(z)$ and OPE:
$$\alpha(z)\alpha(w) \sim \frac{k}{(z-w)^2}$$
 
$$\bar{B}(\mathcal{H}_k) \simeq \text{CE}^!(\mathfrak{h}_k) \quad \text{and} \quad 
  \Omega(\bar{B}(\mathcal{H}_k)) \simeq \mathcal{H}_k$$

where $\text{CE}(\mathfrak{h}_k)$ is the \emph{Chevalley--Eilenberg DG chiral algebra} of the Heisenberg Lie$^*$ algebra.

Thus $(\mathcal{H}_k, \text{CE}(\mathfrak{h}_{-k}))$ form a Koszul pair. The level $k$
parameterizes the central extension, and appears with \emph{inverted sign} as the curvature in the CE algebra (since $h^\vee = 0$ for the Heisenberg):
$$\text{CE}(\mathfrak{h}_{-k}) = V^{\text{CE}}(\mathfrak{h}_{-k}) = (\text{Sym}((s^{-1}N^\vee)_D), d_{\text{CE}}, m_0 = -k \cdot c)$$

The distinction from free fermions is as follows:
\begin{itemize}
\item Free fermions: Simple pole $\psi(z)\psi(w) \sim \frac{1}{z-w}$ → Bar differential is identity → Koszul dual is genuinely different algebra
\item Heisenberg: Double pole $J(z)J(w) \sim \frac{k}{(z-w)^2}$ → Bar differential vanishes → Koszul dual has CE cooperad structure
\end{itemize}

The double pole means:
$$\text{Res}_{z_1=z_2}\left[\frac{k \, dz}{(z_1-z_2)^3}\right] = 0$$
so the bar complex has zero differential except for the curvature term.

Multiple duality structures coexist for the Heisenberg algebra:
\begin{enumerate}
\item \emph{Bar-cobar Koszul duality}: $\mathcal{H}_k^! \simeq \mathrm{Sym}^{\mathrm{ch}}(V^*)$ with curvature $m_0 = -k\omega$
      (equivalently $\mathrm{CE}(\mathfrak{h}_{-k})$ since $\mathfrak{h}$ is abelian; level inverts since $h^\vee = 0$)
\item \emph{Quadratic projection}: $(qP^\circ)^\perp$ gives
      $\text{Sym}((s^{-1}N^\vee)_D)$
      (the underlying graded algebra, missing differential and curvature)
\item \emph{Level-shifting}: $k \leftrightarrow -k$ in representation
      categories (representation theory---different from Koszul duality)
\item \emph{Boson-fermion correspondence}: $V_{\mathbb{Z}} \simeq \mathcal{F}_{\text{charged}}$
      (the lattice VOA $V_{\mathbb{Z}} \supset \mathcal{H}_k$ is isomorphic to the charged free fermion—this is a lattice extension, not Koszul duality)
\end{enumerate}
These are distinct structures; only (1) is bar-cobar Koszul duality.

\begin{remark}[CE versus Sym]
The quadratic projection $(qP^\circ)^\perp$ gives only the quadratic part:
$$\text{Sym}((s^{-1}N^\vee)_D)$$
But the full dual datum $P^{\circ\perp}$ (as in GLZ Proposition 6.2) includes:
\begin{itemize}
\item The differential: $d_{\text{CE}}$ (zero for abelian, but structure still DG)
\item The curving: $m_0 = k \cdot c$ (essential for level dependence)
\item The twisted pair structure: $(B, B^\circ, S)$
\end{itemize}
This is precisely the Chevalley--Eilenberg DG chiral algebra $\text{CE}(\mathfrak{h}_k)$.
See \cite{GLZ-2212.11252v1} Proposition 6.2 (page 19).
\end{remark}

% ================================================================
% NEW SECTION: COMPLETE HEISENBERG KOSZUL DUALITY DERIVATION
% ================================================================

\section{Heisenberg Koszul duality from first principles}
\label{sec:heisenberg-koszul-complete}

\subsection{The Heisenberg example}

The Heisenberg chiral algebra is the simplest non-trivial example in chiral algebra theory,
hence its Koszul dual structure forms an
essential building block:
\begin{enumerate}
\item Heisenberg is simultaneously the abelian case of affine Kac--Moody algebras and also the simplest W-algebra (by degeneracy)
\item It is useful to illustrate the key difference between quadratic projection and full chiral Koszul duality
\item It shows how bar-cobar constructions naturally produce Chevalley--Eilenberg algebras
\item It provides the template for understanding all Lie$^*$ algebra enveloping algebras
\end{enumerate}

\subsection{The setup}

\begin{definition}[Heisenberg Lie$^*$ algebra]
Let $X$ be a smooth algebraic curve. The Heisenberg Lie$^*$ algebra is:
$$\mathfrak{h}_\kappa^* = \mathcal{O}_X \oplus \omega_X \cdot \mathbf{c}$$
with bracket:
$$[J(z), J(w)] = \kappa \cdot \partial_w \delta(z-w) \otimes \mathbf{c}$$
where $J \in \mathcal{O}_X$ is the current and $\mathbf{c}$ is central.
\end{definition}

In OPE language:
$$J(z)J(w) = \frac{\kappa}{(z-w)^2} + \text{regular}$$

The \emph{double pole} is the key feature distinguishing Heisenberg from free fermions.

\subsection{Two perspectives on Heisenberg Koszul duality}

We now derive the Koszul dual using all four perspectives:

\subsubsection{Perspective 1: physical intuition}

Consider free boson CFT with action:
$$S = \frac{\kappa}{4\pi} \int |\partial \phi|^2$$

The current is $J = \partial\phi$. To gauge the $U(1)$ symmetry $\phi \mapsto \phi + \epsilon$:
\begin{itemize}
\item Introduce ghosts: $c$ (fermionic), $b$ (fermionic antighost)
\item BRST operator: $Q = \oint c \partial\phi = \oint cJ$
\item BRST cohomology computes $H^*(\mathfrak{u}(1))$ = Lie algebra cohomology
\end{itemize}

The gauged theory is described by the Chevalley--Eilenberg complex:
$$Q: J \mapsto \kappa \cdot \partial c$$

This is precisely $d_{\text{CE}}$ for the abelian Lie algebra $\mathfrak{h}$.

\emph{Physical conclusion.} Koszul dual is CE(h).

\subsubsection{Perspective 2: geometric intuition}

\emph{Degree~1.}
$\bar{B}_1 = \Gamma(C_2(X), J \boxtimes J \otimes \Omega^1_{\log})$.
The differential is
$$d(J(z_1) \otimes J(z_2) \otimes \eta_{12}) = \text{Res}_{z_1=z_2}[J(z_1)J(z_2) \cdot d\log(z_1-z_2)].$$
Using the OPE:
\begin{align}
J(z_1)J(z_2) \cdot d\log(z_1-z_2) &= \frac{\kappa}{(z_1-z_2)^2} \cdot \frac{dz_1 - dz_2}{z_1-z_2} \\
&= \frac{\kappa \, (dz_1 - dz_2)}{(z_1-z_2)^3}
\end{align}
The triple pole has zero residue:
$$\text{Res}_{z_1=z_2}\left[\frac{\kappa \, dz_1}{(z_1-z_2)^3}\right] = 0.$$
It follows that the bar complex has the following cohomology.

\emph{Degree~0.} $\bar{B}_0 = \mathbb{C} \cdot \mathbf{1}$, $d = 0$.

\emph{Degree~1.} $\bar{B}_1 = \text{span}\{J \otimes J \otimes \eta_{12}\}$,
$d(J \otimes J \otimes \eta_{12}) = 0$ as computed above,
so $H^1 = \bar{B}_1$ survives.

\emph{Degree~2.} $\bar{B}_2 = \text{span}\{J^{\otimes 3} \otimes \eta_{12} \wedge \eta_{23}\}$,
$d(J^{\otimes 3} \otimes \eta_{12} \wedge \eta_{23}) = 0$
by the same double-pole argument, so $H^2 = \bar{B}_2$ survives.



$$d: \bar{B}_1 \to \bar{B}_0 \text{ is the zero map}$$

At every degree, the bar differential vanishes: the double-pole OPE,
combined with the logarithmic form $d\log(z_1 - z_2)$, produces a
triple pole whose residue is zero.  The cohomology is:
$$H^*(\bar{B}(\mathcal{H}_\kappa)) \simeq \bar{B}(\mathcal{H}_\kappa)$$
with CE cooperad structure.


This forces the bar complex to have CE cooperad structure.

\emph{Geometric conclusion.} Bar complex cohomology is $CE(h)$.





\begin{theorem}[Heisenberg Koszul duality; \ClaimStatusProvedElsewhere]
Let $\mathcal{H}_\kappa$ be the Heisenberg chiral algebra at level $\kappa$ on a smooth curve $X$. Then:
$$\mathcal{H}_\kappa^! \simeq \text{CE}(\mathfrak{h}_{-\kappa}) = V^{\text{CE}}(\mathfrak{h}_{-\kappa})$$
where $\text{CE}(\mathfrak{h}_{-\kappa})$ is the Chevalley--Eilenberg DG chiral algebra with:
\begin{itemize}
\item \emph{Underlying space}: $\text{Sym}((s^{-1}N^\vee)_D)$ as graded algebra
\item \emph{Differential}: $d_{\text{CE}} = 0$ (abelian case)
\item \emph{Curvature}: $m_0 = -\kappa \cdot c$ (level inversion: $\kappa \to -\kappa$ for $h^\vee = 0$)
\item \emph{Structure}: CE cooperad in the bar-cobar adjunction
\end{itemize}


\end{theorem}


\subsection{Generalization}

For any Lie$^*$ algebra $\mathfrak{g}$:
$$U(\mathfrak{g})^\kappa \quad \xleftrightarrow{\text{Koszul}} \quad \text{CE}(\mathfrak{g}_{-\kappa-2h^\vee})$$

Heisenberg is the abelian case where $h^\vee = 0$ and $d_{\text{CE}} = 0$, but the CE structure remains.


\subsection{Mathematical significance}

\begin{enumerate}
\item \emph{Lie algebra cohomology}: Chiral algebras naturally encode Lie cohomology via Koszul duality
\item \emph{Deformation theory}: CE algebras control deformations of enveloping algebras
\item \emph{D-modules}: Connection to D-module theory of Lie algebroid actions
\item \emph{Factorization algebras}: CE structure arises naturally from factorization
\end{enumerate}

% ================================================================
% END OF COMPLETE DERIVATION SECTION
% ================================================================


\subsection{Precise definition of chiral Koszul pairs}
\label{subsec:chiral-koszul-pairs-precise}

We now give the definitive definition that applies to all chiral algebras, not just 
quadratic ones.

\begin{definition}[Chiral Koszul pair—Version I: Bar-cobar characterization]
\label{def:chiral-koszul-pair-bar-cobar}
Two chiral algebras $(\mathcal{A}_1, \mathcal{A}_2)$ on a smooth curve $X$ form a 
\emph{chiral Koszul pair} if they satisfy:

\begin{enumerate}
\item \emph{Bar produces dual coalgebra}:
      $$\bar{B}^{\text{ch}}(\mathcal{A}_1) \simeq \mathcal{A}_2^!$$
      as a quasi-isomorphism of chiral coalgebras, where $\mathcal{A}_2^!$ is the 
      Koszul dual coalgebra to $\mathcal{A}_2$

\item \emph{Symmetry}:
      $$\bar{B}^{\text{ch}}(\mathcal{A}_2) \simeq \mathcal{A}_1^!$$
      as a quasi-isomorphism of chiral coalgebras

\item \emph{Cobar reconstructs partner}:
      $$\Omega^{\text{ch}}(\mathcal{A}_2^!) \simeq \mathcal{A}_1 \quad \text{and} \quad
        \Omega^{\text{ch}}(\mathcal{A}_1^!) \simeq \mathcal{A}_2$$
      as quasi-isomorphisms of chiral algebras
\end{enumerate}
\end{definition}

\begin{remark}
\label{rem:why-definition-works}
This definition:
\begin{itemize}
\item \emph{Escapes quadratic constraint}: Makes no reference to presentations by 
      generators and relations
\item \emph{Captures essential duality}: The bar of one is the coalgebra dual to the other
\item \emph{Is geometrically computable}: Configuration spaces provide explicit realizations
\item \emph{Includes classical cases}: Quadratic Koszul pairs satisfy these conditions
\item \emph{Extends to physics}: Natural for vertex operator algebras and CFT
\end{itemize}
\end{remark}

\begin{definition}[Chiral Koszul pair—Version II: Twisting morphism characterization]
\label{def:chiral-koszul-pair-twisting}
\index{twisting morphism|textbf}
Equivalently, $(\mathcal{A}_1, \mathcal{A}_2)$ form a chiral Koszul pair if there exists 
a \emph{universal twisting morphism} $\tau_{12}: \mathcal{A}_1^! \to \mathcal{A}_2$ 
satisfying the Maurer--Cartan equation:
$$d\tau_{12} + \frac{1}{2}[\tau_{12}, \tau_{12}] = 0$$

which induces quasi-isomorphisms:
\begin{align}
\mathcal{A}_1 &\simeq \Omega^{\text{ch}}(\mathcal{A}_2^!)_{\tau_{12}} \\
\mathcal{A}_2 &\simeq (\mathcal{A}_1)_{\tau_{12}}
\end{align}
where subscript $\tau$ denotes twisting by $\tau$.
\end{definition}

\begin{remark}[Twisting morphisms]
\label{rem:twisting-perspective}
\index{Maurer--Cartan element|textbf}
The twisting morphism $\tau_{12}$ is the \emph{explicit map} realizing the Koszul duality:
\begin{itemize}
\item \emph{Domain and codomain}: $\tau_{12}: \mathcal{A}_1^! \to \mathcal{A}_2$ goes 
      from the coalgebra dual to algebra
\item \emph{Maurer--Cartan equation}: Ensures $\tau$ intertwines structures correctly
\item \emph{Geometric realization.}
      $$\tau(c \otimes d) = \int_{\overline{C}_2(X)} \text{ev}^*(c \otimes d) \wedge K(z_1, z_2)$$
      where $K$ is a universal integration kernel
\item \emph{Universality}: Any other twisting factors through $\tau_{12}$
\end{itemize}

This perspective connects to Gui--Li--Zeng's framework where Koszul
duality is expressed through Maurer--Cartan elements in
$\mathcal{A}_1^! \otimes \mathcal{A}_2$.
\end{remark}

\subsection{The Gui--Li--Zeng quadratic duality framework}
\label{subsec:gui-li-zeng-framework}

Our geometric approach to chiral Koszul duality is closely related to the algebraic
framework of Gui, Li, and Zeng \cite{GLZ22}.

\begin{definition}[Gui--Li--Zeng setup]
\label{def:glz}
Gui--Li--Zeng \cite{GLZ22} define Koszul duality for chiral algebras through four ingredients.

\emph{Chiral quadratic data.}
A pair $(N, P)$ where $N$ is a sheaf of generators (chiral vector space) and $P \subset j_* j^* (N \boxtimes N)$ is a subsheaf of quadratic relations.
The quadratic chiral algebra is
$$\mathcal{A}(N, P) = \frac{\mathcal{A}(N)}{(P)}$$
where $\mathcal{A}(N)$ is the free chiral algebra on $N$.

\emph{Dualizable quadratic data.}
$(N, P)$ is \emph{dualizable} if $(s^{-1}N^{\vee}\omega^{-1}, P^{\perp})$ is also a chiral
quadratic datum, where $N^{\vee}\omega^{-1}$ is the dual with twist by inverse canonical bundle and $P^{\perp}$ is the chiral annihilator defined by
$$\mu(\langle P \otimes \omega_{X^2}, P^{\perp} \otimes s^2\omega_{X^2}\rangle) = 0$$
under the unit chiral operation $\mu$.

\emph{Quadratic dual.}
For dualizable $(N, P)$, the quadratic dual is
$$\mathcal{A}^! = \mathcal{A}(s^{-1}N^{\vee}\omega^{-1}, P^{\perp}).$$

\emph{Maurer--Cartan correspondence.}
There is a bijection
$$\operatorname{Hom}(\mathcal{A}, \mathcal{B}) \leftrightarrow MC(\mathcal{A}^! \otimes \mathcal{B})$$
between morphisms of chiral algebras and Maurer--Cartan elements in the tensor product.
\end{definition}

\begin{proposition}[Comparison: our approach vs GLZ; \ClaimStatusProvedHere]
\label{thm:comparison-our-glz}
\index{GLZ framework!comparison}
The geometric bar-cobar framework of this monograph and the algebraic
framework of Gui--Li--Zeng~\cite{GLZ21} are related as follows:
\begin{enumerate}
\item \emph{Quadratic case agreement.}
   For a quadratic chiral algebra $\mathcal{A}(N,P)$ (generators~$N$,
   relations $P \subset j_*j^*(N \boxtimes N)$), our bar construction
   gives $\bar{B}^{\mathrm{ch}}(\mathcal{A}) \simeq
   \mathcal{A}(s^{-1}N^{\vee}\omega^{-1}, P^{\perp})^!$,
   which coincides with the GLZ dual coalgebra.
\item \emph{Non-quadratic extension.}
   Our framework extends to non-quadratic algebras via the full
   configuration space bar complex (Chapter~\ref{chap:bar-cobar}),
   which replaces: quadratic relations $P$ with OPE structure of
   arbitrary pole order; the annihilator $P^{\perp}$ with residue
   extraction at collision divisors; algebraic dualization with
   geometric Poincar\'e--Verdier duality.
\item \emph{Maurer--Cartan elements.}
   The GLZ Maurer--Cartan element
   $\alpha \in MC(\mathcal{A}^! \otimes \mathcal{B})$
   corresponds to our twisting morphism
   $\tau\colon \mathcal{A}^! \to \mathcal{B}$,
   realized geometrically as an integration kernel on
   $\overline{C}_2(X)$:
   $\tau(c)(z) = \int_{\overline{C}_2(X)}
   \mathrm{ev}^* c(w) \wedge K(z, w)$.
\item \emph{Curved structures.}
   \index{$A_\infty$-algebra|textbf}
   GLZ's framework handles curved $A_{\infty}$ algebras through
   Maurer--Cartan deformations; our configuration space approach
   realizes the same data as genus-by-genus quantum corrections
   (curvature $=$ higher-genus corrections; Maurer--Cartan equation
   $=$ Stokes' theorem on $\overline{C}_n(X)$).
\end{enumerate}
\end{proposition}

\begin{proof}
For~(i): the bar complex $\bar{B}^{\mathrm{ch}}(\mathcal{A}(N,P))$
has generators $s^{-1}N^\vee$ in bar degree~1 and relations
$P^\perp \subset s^{-1}N^\vee \otimes s^{-1}N^\vee$ in bar degree~2,
extracted by the double-residue operation on $\overline{C}_3(X)$.  The
$\omega^{-1}$-twist arises from the desuspension convention
(Appendix~\ref{app:sign-conventions}).  Since the algebra is quadratic,
the bar complex is generated in degrees $\leq 2$, so the bar
coalgebra is the sub-coalgebra of the cofree coalgebra
$T^c(s^{-1}N^\vee \omega^{-1})$ cogenerated by $P^\perp$---which is
the definition of the GLZ dual coalgebra
(\cite{GLZ21}, Proposition~3.2).  For~(iii): the MC equation
$d\alpha + \alpha \star \alpha = 0$ translates under the integration
kernel to $d\tau + \mu_{\mathcal{B}} \circ (\tau \otimes \tau)
\circ \Delta_{\mathcal{A}^!} = 0$, which is the twisting morphism
condition (Definition~\ref{def:twisting-morphism}).  Items~(ii)
and~(iv) are descriptions of the respective frameworks' scopes.
\end{proof}

\begin{remark}[Advantages of each approach]
\label{rem:advantages-comparison}

\emph{GLZ algebraic approach:}
\begin{itemize}
\item[+] Clean algebraic formulation
\item[+] Direct definition of dual via annihilators
\item[+] Natural connection to deformation theory
\item[+] Explicit in quadratic case
\item[−] Limited to quadratic or near-quadratic examples
\item[−] Abstract, not immediately computable for complicated algebras
\end{itemize}

\emph{Our geometric approach:}
\begin{itemize}
\item[+] Applies to arbitrary pole order (non-quadratic)
\item[+] Explicitly computable via configuration spaces
\item[+] Natural genus expansion and quantum corrections
\item[+] Physical interpretation via Feynman diagrams
\item[+] Connects to Poincaré--Verdier duality
\item[−] Technically more involved (compactifications, stratifications, Arnold relations)
\item[−] Requires careful analysis of convergence and regularization
\end{itemize}

The two approaches are complementary. GLZ provides conceptual clarity
and algebraic foundations. Our geometric framework provides
computational power and extends to non-quadratic examples essential
for physics (Virasoro, $W$-algebras, Yangian).
\end{remark}

\chapter{Operadic foundations and bar constructions}

 
\section{Symmetric sequences and operads}

\begin{definition}[Symmetric monoidal category]
We work in the symmetric monoidal $\infty$-category $\mathcal{V} = \text{Ch}_\mathbb{C}$ of 
cochain complexes over $\mathbb{C}$ with cohomological grading. The monoidal structure is given by:
\begin{itemize}
\item Unit object: $\mathbb{C}$ concentrated in degree 0
\item Tensor product: $(V \otimes W)^n = \bigoplus_{i+j=n} V^i \otimes W^j$
\item Differential: $d(v \otimes w) = dv \otimes w + (-1)^{|v|}v \otimes dw$
\item Symmetry: $\tau(v \otimes w) = (-1)^{|v||w|}w \otimes v$
\end{itemize}
\emph{Convention.} We use cohomological grading throughout: $\deg(d) = +1$.

All constructions respect this grading and differential structure. For a morphism $f: V \to W$ of degree $|f|$, the Koszul sign rule gives $f(v \otimes w) = (-1)^{|f||v|}f(v) \otimes w$ when extended to tensor products.

% Added for clarity
\emph{Explicit grading convention.} Throughout this paper, we use cohomological grading with $\deg(d) = +1$, and all degree shifts should be interpreted in this context. For a complex $(C^\bullet, d)$, we have $d: C^n \to C^{n+1}$.

\emph{Sign convention for composition.} When composing morphisms of degree $p$ and $q$, the Koszul sign rule applies: passing an element of degree $p$ past one of degree $q$ introduces the sign $(-1)^{pq}$.

\emph{Differential graded context.} All categories are enriched over the category of cochain complexes, with morphism spaces carrying natural differential structures compatible with composition.

\end{definition}

Let $\cV$ be a symmetric monoidal $\infty$-category. In practice, we primarily work with the category of chain complexes over $\C$ (the field of complex numbers), but the constructions apply more generally to any stable presentable symmetric monoidal category. The choice of characteristic 0 is essential for our residue calculus and will be assumed throughout unless otherwise stated.
 
\begin{definition}[Symmetric sequence]
A \emph{symmetric sequence} is a collection $P = \{P(n)\}_{n \geq 0}$ where each $P(n)$ is an object of $\cV$ equipped with a right action of the symmetric group $S_n$. Morphisms of symmetric sequences are collections of $S_n$-equivariant maps. When $\cV$ carries a differential structure, we require that the $S_n$-action commutes with differentials.
\end{definition}
 
The fundamental operation on symmetric sequences is the composition product, which encodes the substitution of operations:
 
\begin{definition}[Composition product]
For symmetric sequences $A$ and $B$, their composition product is defined by:
\[
(A \circ B)(n) = \bigoplus_{k \geq 0} A(k) \otimes_{S_k} \left( \bigoplus_{i_1 + \cdots + i_k = n} \Ind_{S_{i_1} \times \cdots \times S_{i_k}}^{S_n}(B(i_1) \otimes \cdots \otimes B(i_k)) \right)
\]
where $\Ind$ denotes the induced representation functor, using the block diagonal embedding 
\[
S_{i_1} \times \cdots \times S_{i_k} \hookrightarrow S_n
\]
that acts on $\{1, \ldots, i_1\} \sqcup \{i_1 + 1, \ldots, i_1 + i_2\} \sqcup \cdots \sqcup \{i_1 + \cdots + i_{k-1} + 1, \ldots, n\}$.
\end{definition}
 
The composition product is associative up to canonical isomorphism, with unit given by the symmetric sequence $\mathbb{I}$ with $\mathbb{I}(1) = \C$ and $\mathbb{I}(n) = 0$ for $n \neq 1$.
 
\section{Chiral algebras and non-abelian Poincaré duality}
\label{sec:chiral-NAP}

\subsection{Factorization as local-to-global}

\begin{remark}\label{princ:factorization-locality}
The factorization axiom for chiral algebras,
$$\mathcal{A}(U \sqcup V) \simeq \mathcal{A}(U) \otimes \mathcal{A}(V),$$
is the algebraic encoding of locality. From the NAP perspective, this is excision:
$$\int_{M_1 \sqcup M_2} A = \int_{M_1} A \otimes \int_{M_2} A.$$
Factorization algebras are precisely the coefficient systems for non-abelian Poincar\'e duality.
\end{remark}

\subsection{Ran space and universal recipients}

\begin{definition}[Ran space]\label{def:ran-space}
The Ran space of X is:
$$\text{Ran}(X) = \varinjlim_{I \in \text{fSet}^{\text{surj}}} X^I$$
the space of finite non-empty subsets of $X$, where the colimit is taken over the category of non-empty finite sets with surjections. The surjections $I \twoheadrightarrow J$ induce diagonal embeddings $X^J \hookrightarrow X^I$, identifying configurations where points collide.

A factorization algebra on X is equivalent to a constructible sheaf on $\text{Ran}(X)$ satisfying compatibility conditions.
\end{definition}

\begin{theorem}[Chiral algebras on Ran space; \ClaimStatusProvedElsewhere]\label{thm:chiral-ran}
\textup{(Beilinson--Drinfeld, Chapter 3)}

A chiral algebra $\mathcal{A}$ on a curve X determines a D-module $\mathcal{F}_{\mathcal{A}}$ on $\text{Ran}(X)$ satisfying:
\begin{enumerate}
\item Factorization: $\mathcal{F}_{\mathcal{A}}(S \sqcup T) = \mathcal{F}_{\mathcal{A}}(S) \otimes \mathcal{F}_{\mathcal{A}}(T)$
\item Compatibility with embeddings
\item D-module structure encoding OPEs
\end{enumerate}

The bar-cobar duality acts on these sheaves on Ran space, realizing NAP duality for factorization algebras.
\end{theorem}

\begin{remark}[Imported-proof dependency tags]
\emph{Dependency tags.}
\begin{enumerate}[label=(D\arabic*)]
\item Ran-space setup is fixed by Definition~\ref{def:ran-space}.
\item Local-to-global factorization input is Principle~\ref{princ:factorization-locality}.
\item External foundational source is \cite[Chapter~3]{BD04}, where the
chiral-algebra-to-Ran-space package is constructed.
\item Compatibility with this manuscript's duality pipeline is through the
bar-cobar quasi-isomorphism package (Theorem~\ref{thm:bar-cobar-inversion-qi}).
\end{enumerate}
\emph{Provenance note.}
This claim is used as an external theorem, proved in
\cite[Chapter~3]{BD04}, and not reproved here.
\end{remark}

\subsection{Connection to our construction}

Our geometric bar-cobar construction via configuration spaces is the explicit realization of NAP duality for chiral algebras viewed as factorization algebras on Ran space. The configuration space $C_n(X)$ is the n-stratum of Ran(X), and our logarithmic forms encode the factorization structure.

\begin{definition}[Operad]
\index{operad|textbf}
An \emph{operad} $P$ is a monoid for the composition product, equipped with:
\begin{itemize}
\item Composition maps $\gamma : P(k) \otimes P(i_1) \otimes \cdots \otimes P(i_k) \to P(i_1 + \cdots + i_k)$
\item Unit $\eta : \mathbb{I} \to P(1)$ 
\item Associativity axioms ensuring that multi-level compositions are independent of bracketing
\item Equivariance axioms ensuring compatibility with symmetric group actions
\end{itemize}
When $\cV$ has a differential structure, all structure maps must be chain maps.
\end{definition}
 
\begin{definition}[Cooperad]
\index{cooperad|textbf}
A \emph{cooperad} is a comonoid for the composition product, with structure maps dual to those of an operad. Explicitly, we have decomposition maps $\Delta : C(n) \to (C \circ C)(n)$ and a counit $\epsilon : C \to \mathbb{I}$ satisfying coassociativity and coequivariance axioms.
\end{definition}
 
\begin{example}[Endomorphism operad]
For any object $V \in \cV$, the endomorphism operad $\End_V$ has 
\[
\End_V(n) = \Hom_\cV(V^{\otimes n}, V)
\]
with composition given by substitution of multilinear operations. This is the fundamental example motivating the general theory.
\end{example}
 
\section{The cotriple bar construction}
 
Given an adjunction $F \dashv U : \mathcal{A} \rightleftarrows \mathcal{B}$ (with $F$ left adjoint to $U$), we obtain a comonad (also called a cotriple) $G = FU$ on $\mathcal{B}$ with counit $\epsilon : FU \to \text{id}$ and comultiplication $\delta : FU \to FUFU$ induced by the unit and counit of the adjunction.
 
\begin{definition}[Cotriple bar resolution]
The cotriple bar resolution of $B \in \mathcal{B}$ is the simplicial object:
\[
B^G_\bullet(B) : \cdots \rightrightarrows (FU)^3B \rightrightarrows (FU)^2B \rightrightarrows FUB \to B
\]
with face maps $d_i : B^G_n \to B^G_{n-1}$ given by:
\begin{itemize}
\item $d_0 = \epsilon \cdot (FU)^{n-1}$ (apply counit at the first position)
\item $d_i = (FU)^{i-1} \cdot \epsilon \cdot (FU)^{n-i-1}$ for $0 < i < n$ (apply counit at position $i$, collapsing adjacent $FU$'s)
\item $d_n = (FU)^{n-1} \cdot \epsilon$ (apply counit at the last position)
\end{itemize}
and degeneracy maps $s_i : B^G_n \to B^G_{n+1}$ given by inserting the unit of the adjunction at position $i$.
\end{definition}
 
\begin{example}[Operadic bar construction]
For an operad $P$, the free-forgetful adjunction $F_P \dashv U : P\text{-Alg} \rightleftarrows \cV$ yields the classical bar construction $\barB^P_\bullet(A)$ for any $P$-algebra $A$. Explicitly:
\[
\barB^P_n(A) = P \circ \cdots \circ P \circ A \quad \text{($n$ copies of $P$)}
\]
This agrees with the construction via iterated insertions of operations from $P$. The differential is the alternating sum of face maps.
\end{example}
 
\subsection{The fundamental bar-cobar isomorphism}
\index{cobar construction!algebraic|textbf}

The following characterizes Koszul pairs in the classical setting; the chiral analog is Definition~\ref{def:chiral-koszul-pair}.

\begin{remark}
Two objects form a Koszul pair when their bar and cobar constructions are inverses up to quasi-isomorphism: $\barB$ converts algebra structure to coalgebra structure, $\Omega$ converts coalgebra structure to algebra structure, and for a Koszul pair $(A_1, A_2)$ the coalgebra $\barB(A_1)$ is quasi-isomorphic to the coalgebra whose cobar construction recovers $A_2$.
\end{remark}

\begin{definition}[Classical Koszul pair]\label{def:classical-koszul-pair}
\index{homotopy transfer theorem}
Two quadratic operads/algebras $(P_1, P_2)$ with presentations:
\begin{align*}
P_1 &= \mathcal{F}(V_1)/(R_1) \\
P_2 &= \mathcal{F}(V_2)/(R_2)
\end{align*}
form a \emph{Koszul pair} if there exists a perfect pairing $\langle \cdot, \cdot \rangle: V_1 \otimes V_2 \to \mathbb{k}$ such that:

\begin{enumerate}
\item \emph{Generator duality}: $V_2 \cong V_1^* := \text{Hom}(V_1, \mathbb{k})$ via the pairing

\item \emph{Relation orthogonality}: $R_1 \perp R_2$ under the induced pairing on relations

\item \emph{Bar-cobar isomorphism}: There exist quasi-isomorphisms of cooperads and operads:
\begin{align*}
\barB(P_1) &\simeq P_2^! \quad \text{(as cooperads)} \\
\barB(P_2) &\simeq P_1^! \quad \text{(as cooperads)} \\
\Omega(P_1^!) &\simeq P_1 \quad \text{(as operads)} \\
\Omega(P_2^!) &\simeq P_2 \quad \text{(as operads)}
\end{align*}
where $P_i^!$ is the Koszul dual cooperad: the maximal sub-cooperad of $\mathcal{F}^c(V_i^*)$ cogenerated by $R_i^\perp$.
\end{enumerate}
\end{definition}

\begin{remark}
The third condition is the essential content of being a Koszul pair:
the bar construction of $P_1$ computes the dual cooperad $P_2^{!}$
that determines $P_2$. Conversely, $\Omega(B(P_1)) \simeq P_1$.
\end{remark}

\begin{example}[Com-Lie: the prototypical Koszul pair]
For the commutative and Lie operads:
\begin{itemize}
\item Generators: $\mu \in \text{Com}(2)$ (commutative product) and $\ell \in \text{Lie}(2)$ (Lie bracket)
\item Pairing: $\langle \mu, \ell \rangle = 1$ (canonical pairing between symmetry and antisymmetry)
\item Bar-cobar isomorphisms:
\begin{align*}
\barB(\text{Com}) &\simeq \text{co}\text{Lie} \quad \text{(partition complex gives Lie cooperad)} \\
\barB(\text{Lie}) &\simeq \text{co}\text{Com} \quad \text{(Chevalley--Eilenberg gives Com cooperad)} \\
\Omega(\text{co}\text{Lie}) &\simeq \text{Com} \quad \text{(cobar reconstructs commutative structure)} \\
\Omega(\text{co}\text{Com}) &\simeq \text{Lie} \quad \text{(cobar reconstructs Lie structure)}
\end{align*}
\end{itemize}

Concretely: the bar complex of the commutative operad is the chain complex of the partition lattice, whose homology is precisely the Lie operad (with sign).
\end{example}

\begin{remark}
In the chiral setting, we generalize this by:
\begin{itemize}
\item Replacing operads with chiral algebras (factorization algebras on curves)
\item Replacing abstract cooperads with geometric coalgebras (residues on configuration spaces)
\item The isomorphism $\barB(\mathcal{A}_1) \simeq \mathcal{A}_2^!$ becomes a geometric statement about how logarithmic forms (bar side) relate to distributional kernels (cobar side)
\end{itemize}
The principle is the same: Koszul pairs are characterized by bar-cobar being mutually inverse operations.
\end{remark}

\section{The operadic bar-cobar duality}
 
For an augmented operad $P$ with augmentation $\epsilon : P \to \mathbb{I}$, we construct the bar and cobar functors:
 
\begin{definition}[Operadic bar construction]
\index{bar construction!algebraic|textbf}
\index{suspension}
\index{cofree coalgebra}
\index{conilpotent!filtration|textbf}
The bar construction $\barB(P)$ is the cofree cooperad on the desuspension $s^{-1}\bar{P}$ (where $\bar{P} = \ker(\epsilon)$ is the augmentation ideal), in the cohomological convention ($|d| = +1$).  Explicitly:
\[
\barB(P) = T^c(s^{-1}\bar{P}) = \bigoplus_{n \geq 0} (s^{-1}\bar{P})^{\circ n}
\]
where $T^c$ denotes the cofree cooperad functor, $(-)^{\circ n}$ denotes the $n$-fold cooperadic composition,
and the differential $d: \barB(P) \to \barB(P)$ is given by:
\[
d = d_{\text{internal}} + d_{\text{decomposition}}
\]
where:
\begin{itemize}
\item $d_{\text{internal}}$ uses the internal differential of $P$
\item $d_{\text{decomposition}}$ encodes edge contractions on trees decorated with operations from $P$
\end{itemize}
\end{definition}

\section{Cotriple resolution and configuration spaces}

\begin{remark}[Configuration spaces and factorization]
Configuration spaces appear in the bar complex as a consequence of the fundamental theorem of factorization homology (Ayala--Francis \cite{AF}):

\begin{quote}
\emph{``For a factorization algebra $\mathcal{F}$ on a manifold $M$, its factorization homology 
$\int_M \mathcal{F}$ is computed by a Čech-type complex over the Ran space of $M$.''}
\end{quote}

For chiral algebras (2d factorization algebras with conformal structure), this becomes:
$$\int_X \mathcal{A} \simeq \text{colim}_{n} \left[ \mathcal{A}^{\otimes n} \otimes \Omega^*(\text{Conf}_n(X)) \right]$$

The bar complex is precisely the dual construction, explaining its geometric nature.
\end{remark}

\subsection{The genus expansion: a physical and geometric view}
\label{sec:genus_expansion_panorama}

Consider a chiral algebra $\mathcal{A}$ on a curve $X$. The bar-cobar complex
$C_{\bullet}(\mathcal{A})$ involves tensor products of $\mathcal{A}$ at distinct
points of $X$:
$$\mathcal{A}_{x_1} \otimes \mathcal{A}_{x_2} \otimes \cdots \otimes \mathcal{A}_{x_n}$$
The genus of the underlying surface determines the relevant correction terms:

\begin{itemize}
\item \emph{Genus 0 (Tree level):} Points connected by a sphere --- this gives
the classical bar complex, the associative structure.

\item \emph{Genus 1 (One loop):} Points connected by a torus --- this is where
\emph{central extensions} first appear. The trace $\operatorname{Tr}(a \otimes b)$
around the $S^1$ of the torus encodes the central charge.

\item \emph{Genus $g \geq 2$ (Multiple loops):} Surfaces with multiple handles ---
higher genus corrections to the OPE, encoding modular structure.
\end{itemize}

\subsubsection{The geometric construction}

Consider the configuration spaces:
$$\mathrm{Conf}_n(\Sigma_g) = \{ (x_1, \ldots, x_n) \in \Sigma_g^n \mid x_i \neq x_j \text{ for } i \neq j \}$$
for $\Sigma_g$ a Riemann surface of genus $g$.

The \emph{genus $g$ bar complex} is precisely:
$$C_{\bullet}^{(g)}(\mathcal{A}) = \int_{\mathrm{Conf}_{\bullet}(\Sigma_g)} 
\mathcal{A}^{\boxtimes \bullet}$$
where the integration is factorization homology in the sense of Ayala--Francis.

\subsubsection{The functorial uniqueness}

The genus stratification is not a choice but a necessity.
The category of chiral algebras naturally extends to a category of \emph{modular
chiral algebras}, where operations are parametrized by:
$$\mathcal{P}(g,n) = \text{moduli of genus-}g\text{ curves with }n\text{ marked points}$$

The functor:
$$\mathcal{A} \mapsto \{ C_{\bullet}^{(g)}(\mathcal{A}) \}_{g \geq 0}$$
is uniquely determined by:
\begin{enumerate}
\item Functoriality under degenerations $\Sigma_g \rightsquigarrow \Sigma_{g-1}$
(separating a handle)
\item Compatibility with factorization
\item Genus 0 data (the classical structure)
\end{enumerate}

\subsubsection{The physical interpretation}

In conformal field theory, the genus expansion \emph{is} the loop expansion:
$$Z_{\text{CFT}} = \sum_{g=0}^{\infty} \hbar^{2g-2} \int_{\mathcal{M}_g} F_g$$
where $\mathcal{M}_g$ is the moduli space of genus-$g$ curves.

Our bar-cobar construction at genus $g$ computes exactly the integrand $F_g$.
The central charge $\kappa$ plays the role of $\hbar$.

\begin{theorem}[Operadic bar complex; \ClaimStatusProvedElsewhere]\label{thm:operadic-bar}
\index{operadic bar construction}
For an operad $\mathcal{P}$ and $\mathcal{P}$-algebra $A$, the bar complex is:
$$B_{\mathcal{P}}(A) = \bigoplus_{n \geq 1} (\mathcal{P}^{\text{\raisebox{0.5ex}{\tiny !`}}}(n) \otimes_{\Sigma_n} A^{\otimes n})$$
where $\mathcal{P}^{\text{\raisebox{0.5ex}{\tiny !`}}}$ is the Koszul dual cooperad (which already incorporates the appropriate suspension), and the differential combines the cooperad cocomposition with the algebra structure maps of $A$.
\end{theorem}

\begin{theorem}[Geometric realization; \ClaimStatusProvedHere]\label{thm:geometric-bridge}
For the chiral operad $\mathcal{P}_{\text{ch}}$ on a curve $X$:
\begin{enumerate}
\item $\mathcal{P}_{\text{ch}}(n) \cong \Omega^*_{\log}(\overline{C}_n(X))$ as a complex (Kontsevich--Soibelman)
\item The operadic composition corresponds to boundary stratification
\item The bar differential becomes residues at collision divisors
\end{enumerate}

This provides a canonical isomorphism:
$$B_{\mathcal{P}_{\text{ch}}}(\mathcal{A}) \cong \bar{B}^{\text{ch}}_{\text{geom}}(\mathcal{A})$$
\end{theorem}

\begin{proof}
We verify the isomorphism in two steps.

\emph{Underlying graded objects.}
The Koszul dual cooperad of the chiral operad satisfies
$\mathcal{P}_{\mathrm{ch}}^{!'}(n)
\cong H_*(\overline{C}_n(X), \mathbb{C})$
(the Borel--Moore homology of the FM compactification, dual to the
logarithmic cohomology
$\mathcal{P}_{\mathrm{ch}}(n) = \Omega^*_{\log}(\overline{C}_n(X))$).
Therefore the $n$-th component of the operadic bar complex is:
\[
\mathcal{P}_{\mathrm{ch}}^{!'}(n) \otimes_{\Sigma_n} \mathcal{A}^{\otimes n}
\cong \Gamma\bigl(\overline{C}_n(X),\,
j_* j^*(\mathcal{A}^{\boxtimes n}) \otimes \Omega^*_{\log}\bigr)^{\Sigma_n}
\]
which is the $n$-th term of the geometric bar complex
$\bar{B}^{\mathrm{ch}}_{\mathrm{geom}}(\mathcal{A})$.

\emph{Differentials.}
The operadic bar differential $d = d_{\mathrm{cooperad}} +
d_{\mathrm{alg}}$ translates as follows.
The cooperad cocomposition
$\Delta\colon \mathcal{P}^{!'}(n)
\to \bigoplus_{k} \mathcal{P}^{!'}(k) \otimes
\mathcal{P}^{!'}(n_1) \otimes \cdots \otimes \mathcal{P}^{!'}(n_k)$
corresponds to the boundary stratification of $\overline{C}_n(X)$:
each codimension-$1$ boundary divisor $D_S$ (labeled by a subset $S$ of
colliding points) contributes a residue term.  The algebra structure maps
$m_k\colon \mathcal{A}^{\otimes k} \to \mathcal{A}$ are the OPE
coefficients extracted by these residues.  Together,
$d_{\mathrm{cooperad}} + d_{\mathrm{alg}}$ becomes
$d_{\mathrm{dR}} + d_{\mathrm{res}}$
(the de~Rham plus residue differential on the geometric bar complex),
matching the differential of
Definition~\ref{def:geometric-bar}.
\end{proof}

\section{Com-Lie Koszul duality from first principles}
 
\section{Quadratic operads and Koszul duality}
 
We now specialize to quadratic operads, which admit a particularly refined duality theory:
 
\begin{definition}[Quadratic operad]
A quadratic operad has the form $P = \Free(E)/(R)$ where:
\begin{itemize}
\item $E$ is a collection of generating operations concentrated in arity 2
\item $R \subset \Free(E)(3)$ consists of quadratic relations (involving exactly two compositions)
\item $\Free$ denotes the free operad functor
\item $(R)$ denotes the operadic ideal generated by $R$
\end{itemize}
\end{definition}
 
\begin{definition}[Koszul dual cooperad]
\index{desuspension}
The Koszul dual cooperad $P^!$ is the maximal sub-cooperad of the cofree cooperad $T^c(s^{-1}E^\vee)$ cogenerated by the orthogonal relations $R^\perp \subset (s^{-1}E^\vee)^{\otimes 2}$, where the orthogonality is with respect to the natural pairing induced by evaluation.
\end{definition}
 
\begin{definition}[Koszul operad]\label{def:koszul-operad}
\index{Koszul operad|textbf}
An operad $P$ is \emph{Koszul} if the canonical map $\Omega(P^!) \to P$ is a quasi-isomorphism. Equivalently, the Koszul complex $K_\bullet(P) = P^! \circ P$ with differential induced by the cooperad and operad structures is acyclic in positive degrees.
\end{definition}
 
\section{Derivation of Com-Lie duality}
 
We now prove the fundamental duality between the commutative and Lie operads:
 
\begin{theorem}[Com-Lie Koszul duality; \ClaimStatusProvedHere]\label{thm:com-lie}
\index{$L_\infty$-algebra}
We have canonical isomorphisms of operads:
\[
\Com^! \cong \Lie \quad \text{and} \quad \Lie^! \cong \Com
\]
At the cooperad level, $\barB(\Com) \simeq \mathrm{co}\Lie$ and $\barB(\Lie) \simeq \mathrm{co}\Com$.
Moreover, both $\Com$ and $\Lie$ are Koszul operads with quasi-isomorphisms:
\[
\Omega(\text{co}\Lie) \xrightarrow{\sim} \Com, \quad \Omega(\text{co}\Com) \xrightarrow{\sim} \Lie
\]
\end{theorem}
 
\begin{proof}[Proof via partition lattices]
By Theorem \ref{thm:partition}, $\barB(\Com)(n) \simeq s^{n-1}\tilde{C}_{*}(\barPi_n) \otimes \sgn_n$.
 
Classical results of Björner--Wachs \cite{BW93} and Stanley \cite{Sta97} establish that the reduced homology of $\barPi_n$ is:
\begin{itemize}
\item The order complex $|\barPi_n|$ has dimension $n-3$ (maximal chains in $\barPi_n$ have length $n-3$)
\item $\tilde{H}_k(\barPi_n) = 0$ for $k \neq n-3$ (homology concentrated in the top degree)
\item The $S_n$-representation on $\tilde{H}_{n-3}(\barPi_n)$ is $\Lie(n) \otimes \sgn_n$ where $\Lie(n)$ is the Lie representation
\end{itemize}

The geometric realization $|\barPi_n|$ is homotopy equivalent to a wedge of $(n-1)!$ spheres of dimension $n-3$, with the $S_n$-action on the top homology given by the Lie representation tensored with the sign.

For Com-Lie duality, the bar construction gives:
$$\barB(\Com)(n) \simeq s^{n-1}\tilde{C}_{*}(\barPi_n) \otimes \sgn_n$$
Taking homology and using that $\barPi_n$ has homology concentrated in the single degree $n-3$:
$$H_*(\barB(\Com)(n)) \simeq s^{n-1}\tilde{H}_{n-3}(\barPi_n) \otimes \sgn_n = s^{n-1}\Lie(n) \otimes \sgn_n \otimes \sgn_n = s^{n-1}\Lie(n)$$
Since this is concentrated in a single degree, the bar complex is formal and we obtain:
$$\barB(\Com) \simeq \text{co}\Lie[1]$$
as required.
 
Since the bar complex has homology concentrated in a single degree, it follows that:
\[
H_*(\barB(\Com)) \cong \text{co}\Lie[1]
\]
where the shift accounts for the suspension. Applying $\Omega$ yields $\Omega(\text{co}\Lie) \simeq \Com$.
 
The dual statement $\Lie^! \cong \text{co}\Com$ follows by Schur--Weyl duality, using the characterization of $\Lie$ as the primitive part of the tensor coalgebra.
\end{proof}
 
\begin{proof}[Alternative proof via generating series]
The Poincaré series of the operads satisfy:
\begin{align}
P_{\Com}(x) &= e^x - 1 \\
P_{\Lie}(x) &= -\log(1 - x)
\end{align}
These are compositional inverses: $P_{\Lie}(-P_{\Com}(-x)) = x$. This functional equation characterizes Koszul dual pairs, providing an independent verification of the duality.
\end{proof}
 
\section{The quadratic dual and orthogonality}
 
For explicit computations, we need the quadratic presentations:
 
\begin{proposition}[Quadratic presentations; \ClaimStatusProvedElsewhere]
The operads $\Com$ and $\Lie$ have quadratic presentations:
\begin{align}
\Com &= \Free(\mu)/(R_{\Com}) \text{ where } R_{\Com} = \langle \mu_{12,3} - \mu_{1,23}, \mu_{12} - \mu_{21} \rangle \\
\Lie &= \Free(\ell)/(R_{\Lie}) \text{ where } R_{\Lie} = \langle \ell_{12,3} + \ell_{23,1} + \ell_{31,2}, \ell_{12} + \ell_{21} \rangle
\end{align}
where subscripts denote inputs, and composition is denoted by adjacency. Here $\mu_{12,3}$ means $\mu \circ_1 \mu$ and $\mu_{1,23}$ means $\mu \circ_2 \mu$.
\end{proposition}
 
\begin{proposition}[Orthogonality; \ClaimStatusProvedHere]\label{prop:orthogonal}
With respect to the natural pairing between $\Free(\mu)(3)$ and
$\Free(\ell^*)(3)$ induced by $\langle \mu, \ell^* \rangle = 1$, one has
\[
R_{\Com} \perp R_{\Lie}
\]
This orthogonality is the concrete manifestation of Koszul duality.
\end{proposition}
 
\begin{proof}
We compute the pairing explicitly. The spaces have bases:
\begin{align}
\Free(\mu)(3) &= \text{span}\{\mu_{12,3}, \mu_{1,23}, \mu_{13,2}, \mu_{2,13}, \mu_{23,1}, \mu_{3,12}\} \\
\Free(\ell^*)(3) &= \text{span}\{\ell^*_{12,3}, \ell^*_{1,23}, \text{etc.}\}
\end{align}
 
The pairing $\langle \mu_{ij,k}, \ell^*_{pq,r} \rangle = 1$ if the tree structures match and $0$ otherwise. Computing:
\begin{align}
\langle \mu_{12,3} - \mu_{1,23}, \ell^*_{12,3} + \ell^*_{23,1} + \ell^*_{31,2} \rangle &= 1 + 0 + 0 - 0 - 0 - 1 = 0 \\
\langle \mu_{12,3} - \mu_{1,23}, \ell^*_{13,2} + \ell^*_{32,1} + \ell^*_{21,3} \rangle &= 0 - 1 + 0 + 0 + 1 + 0 = 0
\end{align}
Similar computations for all pairs verify the orthogonality.
\end{proof}

\section{Factorization algebra axioms: verification}
\label{sec:factorization-axioms-complete}

\subsection{Motivation}

\begin{remark}[Physical locality]
In quantum field theory, locality requires that observables in spacelike separated
regions commute (or anti-commute for fermions).  Mathematically, for disjoint
regions $U, V \subset M$ this becomes
$$\mathcal{F}(U \sqcup V) \cong \mathcal{F}(U) \otimes \mathcal{F}(V),$$
which is the factorization axiom.  Factorization algebras formalize the assembly
of local observables into a global theory.
\end{remark}

\begin{construction}[Configuration space realization]
Factorization algebras are \emph{coefficient systems} on configuration spaces.

For a manifold $M$, consider the configuration space:
$$C_n(M) = \{(x_1, \ldots, x_n) \in M^n : x_i \neq x_j \text{ for } i \neq j\}$$

A factorization algebra $\mathcal{F}$ assigns:
\begin{itemize}
\item To each open $U \subset M$: a vector space $\mathcal{F}(U)$
\item To each configuration $(x_1, \ldots, x_n) \in C_n(U)$: structure maps
\end{itemize}

The factorization property encodes:
$$\mathcal{F}(U)= \text{colim}_{(x_1,\ldots,x_n) \in C_n(U)} 
   \mathcal{F}(\text{disk around } x_1) \otimes \cdots \otimes 
   \mathcal{F}(\text{disk around } x_n)$$

The algebraic structure arises from the geometry of configuration spaces.
\end{construction}

\begin{computation}[Explicit verification; \ClaimStatusProvedHere]
We verify the factorization axioms explicitly for:

\emph{Observables in mechanics.}
$$\mathcal{F}(U) = C^\infty(U) \quad \text{(functions on configuration space)}$$

For disjoint $U, V$:
$$\mathcal{F}(U \sqcup V) = C^\infty(U \sqcup V) \cong C^\infty(U) \oplus C^\infty(V)
   = \mathcal{F}(U) \oplus \mathcal{F}(V)$$
(Note: this is a direct sum/product, not a tensor product. For the tensor factorization axiom $\mathcal{F}(U \sqcup V) \cong \mathcal{F}(U) \otimes \mathcal{F}(V)$, one instead uses the symmetric algebra $\mathcal{F}(U) = \text{Sym}(\Gamma_c(U, E))$ on compactly supported sections.)

\emph{Chiral algebra (Heisenberg).}
$$\mathcal{H}(U) = \text{Free chiral algebra generated by } U$$

For disjoint disks $D_1, D_2 \subset \mathbb{C}$:
$$\mathcal{H}(D_1 \sqcup D_2) = \mathcal{H}(D_1) \otimes \mathcal{H}(D_2)$$

(no interaction between separated regions).
\end{computation}

\begin{remark}
A factorization algebra $\mathcal{F}$ on $M$ is the initial object in the
category of functors $\mathrm{Opens}(M) \to \mathcal{V}$ satisfying locality (factorization for disjoint unions) and compatible with inclusions. This universal property determines factorization algebras uniquely up to
canonical isomorphism. Factorization algebras on $\mathbb{R}^n$ are
equivalent to $E_n$-algebras (algebras over the little $n$-disks operad)~\cite[Theorem~5.4.5.9]{HA}.
\end{remark}

\subsection{Weiss topology and factorization algebras}

The cosheaf condition defining factorization algebras requires a topology
finer than the ordinary open-cover topology.  The correct notion,
introduced by Weiss~\cite{Wei99} in the context of manifold calculus and
adopted by Costello--Gwilliam~\cite{CG17}, is the following.

\begin{definition}[Weiss cover]\label{def:weiss-cover}
\index{Weiss cover|textbf}
Let $M$ be a topological space.  A collection $\{U_i\}_{i \in I}$ of open
subsets of $M$ is a \emph{Weiss cover} of an open $V \subset M$ if every
\emph{finite} subset $S \subset V$ is contained in some single $U_i$.
\end{definition}

\begin{remark}[Weiss covers versus ordinary covers]\label{rem:weiss-vs-ordinary}
Every Weiss cover is an ordinary open cover (take $|S|=1$), but the
converse fails: two disjoint opens $U_1, U_2$ with $U_1 \cup U_2 = V$
form an ordinary cover of~$V$ but not a Weiss cover (any two-point set
$\{p_1, p_2\}$ with $p_i \in U_i \setminus U_j$ is contained in
neither~$U_i$).  The Weiss condition forces the opens to be
``large enough'' to contain arbitrarily many points simultaneously, which
is precisely the condition needed for the factorization product to detect
all multi-point interactions.
\end{remark}

The \emph{Weiss topology} on $\mathrm{Opens}(M)$ is the Grothendieck
topology whose covering families are exactly the Weiss covers.
A \emph{prefactorization algebra} that satisfies descent for the Weiss
topology is a \emph{factorization algebra}.

\subsection{Ayala--Francis axioms}

\begin{definition}[Factorization algebra --- Ayala--Francis definition]
\label{def:factorization-algebra-AF}
Let $M$ be a smooth manifold and $\mathcal{V}$ a symmetric monoidal $\infty$-category
(in the sense of~\cite[§2.0, §4.1]{HA}).
A \emph{factorization algebra} $\mathcal{F}$ on $M$ with values in $\mathcal{V}$
consists of:

\emph{Data.}
\begin{enumerate}
\item For each open $U \subset M$: an object $\mathcal{F}(U) \in \mathcal{V}$.

\item For each inclusion $U \hookrightarrow V$ of opens: a morphism
   $\mathcal{F}(U) \to \mathcal{F}(V)$ in $\mathcal{V}$.

\item For each finite collection of pairwise disjoint opens $U_1, \ldots, U_n \subset V$:
   a \emph{factorization map}
   $$\mu_{U_1,\ldots,U_n}^V: \mathcal{F}(U_1) \otimes \cdots \otimes \mathcal{F}(U_n)
      \to \mathcal{F}(V).$$
\end{enumerate}

\emph{Axioms.}

(FA1) \emph{Functoriality.} The assignment $U \mapsto \mathcal{F}(U)$ is a functor
from $\text{Opens}(M)$ to $\mathcal{V}$:
\begin{itemize}
\item $\mathcal{F}(U) \xrightarrow{\text{id}} \mathcal{F}(U)$ is the identity.
\item For $U \hookrightarrow V \hookrightarrow W$: the composition
   $\mathcal{F}(U) \to \mathcal{F}(V) \to \mathcal{F}(W)$ equals
   $\mathcal{F}(U) \to \mathcal{F}(W)$.
\end{itemize}

(FA2) \emph{Multiplicativity.} For disjoint opens $U_1, \ldots, U_n \subset V$,
the factorization map is a structure map:
$$\mu_{U_1,\ldots,U_n}^V: \mathcal{F}(U_1) \otimes \cdots \otimes \mathcal{F}(U_n)
   \longrightarrow \mathcal{F}(V)$$
This map is a quasi-isomorphism when $\{U_1, \ldots, U_n\}$ forms a
Weiss cover of $V$ (Definition~\ref{def:weiss-cover}), but for arbitrary disjoint opens it is merely a morphism (not necessarily an equivalence).

(FA3) \emph{Associativity.} For nested collections
$U_{ij} \subset V_i \subset W$ (all disjoint), the diagram commutes:
$$\begin{tikzcd}[column sep=large]
\bigotimes_{i,j} \mathcal{F}(U_{ij})
   \arrow[r, "{\otimes_i \mu_{U_{i,*}}^{V_i}}"]
   \arrow[d, "{\mu_{\{U_{ij}\}}^W}"'] &
\bigotimes_i \mathcal{F}(V_i)
   \arrow[d, "{\mu_{\{V_i\}}^W}"] \\
\mathcal{F}(W) \arrow[r, equal] & \mathcal{F}(W)
\end{tikzcd}$$

(FA4) \emph{Unit.} For any open $U$:
$$\mathcal{F}(\emptyset) = \mathbb{1}_{\mathcal{V}}
   \quad \text{(unit object in } \mathcal{V}\text{)}$$
and $\mathcal{F}(\emptyset) \otimes \mathcal{F}(U) \xrightarrow{\sim} \mathcal{F}(U)$.

(FA5) \emph{Symmetry.} For any permutation $\sigma \in S_n$ and opens
$U_1, \ldots, U_n \subset V$, the diagram commutes:
$$\begin{tikzcd}
\mathcal{F}(U_1) \otimes \cdots \otimes \mathcal{F}(U_n)
   \arrow[r, "\sigma"] \arrow[d, "\mu"'] &
\mathcal{F}(U_{\sigma(1)}) \otimes \cdots \otimes \mathcal{F}(U_{\sigma(n)})
   \arrow[d, "\mu"] \\
\mathcal{F}(V) \arrow[r, equal] & \mathcal{F}(V)
\end{tikzcd}$$
\end{definition}

\begin{remark}[Interpretation of axioms]
(FA1) says $\mathcal{F}$ is a presheaf; (FA2) that observables on disjoint regions are independent (locality); (FA3) that the order of combining observables does not matter; (FA4) that the empty region contributes trivially; and (FA5) that the labeling of regions is inessential.
\end{remark}

\subsection{Chiral algebras as factorization algebras}

\begin{theorem}[Chiral algebras are factorization algebras; \ClaimStatusProvedElsewhere]
\label{thm:chiral-factorization}
\textup{(Beilinson--Drinfeld \cite[Theorem~3.4.9, \S3.4.10--3.4.12]{BD04})}

Let $X$ be a smooth algebraic curve over $\mathbb{C}$.
Every chiral algebra $\mathcal{A}$ on $X$ in the sense of
\cite[\S3.3]{BD04} determines a factorization algebra of
$\mathcal{D}$-modules on $\operatorname{Ran}(X)$ satisfying
axioms (FA1)--(FA5).  More precisely, the functor
$$\mathcal{A} \;\longmapsto\; \mathcal{F}_{\mathcal{A}},
\qquad \mathcal{F}_{\mathcal{A}}(U) = \Gamma(U, \mathcal{A}),$$
defines an equivalence between the category of chiral algebras
on $X$ and the category of factorization algebras of
$\mathcal{D}$-modules on $\operatorname{Ran}(X)$.
\end{theorem}

\begin{remark}[Dictionary between BD and the present conventions]
\label{rem:bd-factorization-dictionary}
The equivalence of Theorem~\ref{thm:chiral-factorization} is
\cite[Theorem~3.4.9]{BD04}.  We recall the translation between
Beilinson--Drinfeld's notation and ours:
\begin{center}
\renewcommand{\arraystretch}{1.25}
\begin{tabular}{lll}
\textbf{This text} & \textbf{BD \cite{BD04}} & \textbf{Reference} \\
\hline
Factorization algebra $\mathcal{F}$ & Factorization $\mathcal{D}_X$-algebra &
  \cite[\S3.4.1]{BD04} \\
Multiplicativity (FA2) & Factorization isomorphism &
  \cite[(3.4.4.2)]{BD04} \\
Coherence (FA3) & Higher compatibilities &
  \cite[Proposition~3.4.2]{BD04} \\
$S_n$-equivariance (FA5) & Equivariant structure &
  \cite[\S3.4]{BD04} \\
Chiral product $\mu$ & Chiral operation $\mu^{\operatorname{ch}}$ &
  \cite[\S3.3.1]{BD04} \\
Configuration space $C_n(X)$ & $n$-th stratum of $\operatorname{Ran}(X)$ &
  \cite[\S3.4.10--3.4.12]{BD04} \\
\end{tabular}
\end{center}

The translation proceeds as follows:
\begin{enumerate}
\item BD's \emph{factorization isomorphism}
$\mathcal{A}|_{U_1 \sqcup \cdots \sqcup U_n} \cong
\mathcal{A}|_{U_1} \boxtimes \cdots \boxtimes \mathcal{A}|_{U_n}$
(\cite[(3.4.4.2)]{BD04}) is precisely our axiom (FA2), with
$\boxtimes$ the external tensor product of $\mathcal{D}$-modules
on disjoint opens reducing to the ordinary tensor product on
global sections.
\item Coherence of the factorization isomorphisms
(\cite[Proposition~3.4.2]{BD04}) gives axiom (FA3).
\item The $S_n$-equivariance built into BD's definition
(\cite[\S3.4]{BD04}) gives axiom (FA5).
\item Axioms (FA1) and (FA4) are the presheaf and unit conditions,
which are part of the $\mathcal{D}$-module structure.
\end{enumerate}

The remaining content of this monograph---in particular the
bar-cobar constructions via configuration space integrals on
Fulton--MacPherson compactifications---provides an
\emph{explicit geometric realization} of the factorization
structure, which is not contained in~\cite{BD04}.
\end{remark}

\subsection{Gluing formulas and excision}

\begin{theorem}[Excision property; \ClaimStatusProvedElsewhere]
\label{thm:excision-factorization}
\textup{(Costello--Gwilliam~\cite[§6.3]{CG17}; Lurie~\cite[Theorem~5.5.3.11]{LurieHA})}

Let $\mathcal{F}$ be a factorization algebra on $M$, and let
$M = U \cup V$ be a \emph{Weiss cover} (i.e., every finite subset of $M$ is contained in $U$ or in $V$). Then there is a natural equivalence:
$$\mathcal{F}(M) \simeq \mathcal{F}(U) \otimes_{\mathcal{F}(U \cap V)} \mathcal{F}(V)$$
where the tensor product is the derived pushout over $\mathcal{F}(U \cap V)$.

For \emph{locally constant} factorization algebras ($E_n$-algebras), excision holds more generally for collar-gluing decompositions of manifolds, without the Weiss hypothesis~\cite[Theorem~5.5.3.11]{LurieHA}.
\end{theorem}

\begin{proof}
This is a special case of the cosheaf property for Weiss covers
(Theorem~\ref{thm:factorization-cosheaf} below) applied to the
two-element Weiss cover $\{U, V\}$.  The pushout diagram
$$\begin{tikzcd}
U \cap V \arrow[r] \arrow[d] & V \arrow[d] \\
U \arrow[r] & M
\end{tikzcd}$$
in $\mathrm{Opens}(M)$ is sent by $\mathcal{F}$ to a pushout in the target
category by the descent condition for Weiss covers.
\end{proof}

\subsection{Cosheaf property}

\begin{theorem}[Factorization algebras are cosheaves for Weiss covers;
\ClaimStatusProvedElsewhere]
\label{thm:factorization-cosheaf}
\textup{(Costello--Gwilliam \cite[Theorem~6.3.1]{CG17})}

A prefactorization algebra $\mathcal{F}$ on $M$ is a factorization algebra
if and only if it satisfies the \emph{cosheaf property for Weiss covers}:
for every Weiss cover $\{U_i\}_{i \in I}$ of an open $V \subset M$
(Definition~\ref{def:weiss-cover}), the natural map
$$\operatorname{hocolim}_{\check{C}(\{U_i\})}
   \mathcal{F} \;\xrightarrow{\;\sim\;}\; \mathcal{F}(V)$$
is an equivalence, where the homotopy colimit is taken over the
\v{C}ech nerve of the Weiss cover.
\end{theorem}

\begin{remark}[Weiss covers versus ordinary covers]\label{rem:why-weiss}
The cosheaf condition for \emph{ordinary} open covers is too weak: it
detects only single-point behavior and misses the multi-point interactions
encoded by the factorization product.  The Weiss condition ensures that
every finite configuration of points is witnessed by some single open in
the cover, so that the factorization maps $\mu_{U_1,\ldots,U_n}^V$ are
fully probed.  Conversely, the cosheaf condition for Weiss covers is
\emph{exactly equivalent} to the factorization axioms (FA1)--(FA5)
together with appropriate descent; see~\cite[Proposition~6.3.4]{CG17}.
\end{remark}

\subsection{Master verification table}

\begin{table}[H]
\centering
\caption{Factorization algebra axioms}
\begin{tabular}{|l|c|c|}
\hline
\textbf{Axiom} & \textbf{Statement} & \textbf{Verification for Chiral Algebras} \\
\hline
(FA1) Functoriality & $U \to V$ gives $\mathcal{F}(U) \to \mathcal{F}(V)$ & 
 (restriction maps) \\
\hline
(FA2) Multiplicativity & $\mathcal{F}(U_1 \sqcup \cdots \sqcup U_n) \cong 
\bigotimes_i \mathcal{F}(U_i)$ & (BD factorization) \\
\hline
(FA3) Associativity & Multi-level factorization commutes & (coherence) \\
\hline
(FA4) Unit & $\mathcal{F}(\emptyset) = \mathbb{1}$ & (vacuum vector) \\
\hline
(FA5) Symmetry & Permutation equivariance & ($S_n$-equivariance) \\
\hline
\textbf{Excision} & $\mathcal{F}(U \cup V) = \mathcal{F}(U) \otimes_{\mathcal{F}(U \cap V)} 
\mathcal{F}(V)$ & (Mayer--Vietoris) \\
\hline
\textbf{Cosheaf} & $\text{colim}_I \bigotimes_{i \in I} \mathcal{F}(U_i) \to 
\mathcal{F}(V)$ & (local-to-global) \\
\hline
\end{tabular}
\end{table}

\bigskip
\noindent\textbf{What we have established.}\enspace
The algebraic foundations of Koszul duality are now in place:
the bar-cobar adjunction, the Koszul criterion via diagonal Ext
vanishing, the operadic framework for $\mathrm{Com}^! = \mathrm{Lie}$ and
its relatives, and the factorization axioms that govern local-to-global
assembly.  These are the genus-$0$ ingredients.  In
Chapter~\ref{chap:config-spaces}, the tensor coalgebra $T^c(s^{-1}\bar{A})$
will be replaced by sections of $\cA^{\boxtimes n}$ on
Fulton--MacPherson compactifications $\overline{C}_n(X)$, the
alternating-sum differential will become a sum of residues at collision
divisors, and the Koszul criterion will acquire a geometric proof via
Arnold relations.  The abstract adjunction $\bar{B} \dashv \Omega$
will then become the Verdier duality on $\operatorname{Ran}(X)$
that drives the entire monograph.

\subsection{Summary and significance}

